% Reachability: the GC keeps what the roots can reach; the rest is
% garbage. Used in M10-L02.
\documentclass[tikz,border=4pt]{standalone}
\usepackage{tikz}
\input{_m10-styles.texinput}
\begin{document}
\begin{tikzpicture}[m10, x=1mm, y=-1mm, font=\large]
  % roots
  \node[stackcell, minimum width=30mm, minimum height=26mm, font=\large] (roots) at (0,22)
    {roots\\[3pt]stack,\\registers,\\globals};

  % reachable blocks
  \node[heapcell, minimum width=20mm, minimum height=12mm, font=\large] (a) at (52,6) {a};
  \node[heapcell, minimum width=20mm, minimum height=12mm, font=\large] (b) at (88,6) {b};
  \node[heapcell, minimum width=20mm, minimum height=12mm, font=\large] (c) at (70,26) {c};
  \node[goodnote, anchor=west, font=\large] at (104,10) {reachable = kept};

  % garbage blocks
  \node[freedcell, hatched, minimum width=20mm, minimum height=12mm] (x) at (52,46) {};
  \node[freedcell, hatched, minimum width=20mm, minimum height=12mm] (y) at (88,46) {};
  \node[badnote, anchor=west, font=\large] at (104,46) {unreachable = garbage};

  \draw[flow, line width=1pt] (roots.east) -- (a.west);
  \draw[flow, line width=1pt] (a.east) -- (b.west);
  \draw[flow, line width=1pt] (a.south) -- (c.north);
  \draw[badflow, line width=1pt] (x.east) -- (y.west);

  \node[note, anchor=north west, font=\large\itshape] at (0,60)
    {A reachable pointer keeps its block alive, so no dangling pointer can exist.};
\end{tikzpicture}
\end{document}
